% ================================================================================
% ANTECEDENTES
% ================================================================================

\section{Investigación previa en calificación predictiva de leads}

La sustitución de sistemas de calificación basados en reglas por modelos ajustados
sobre el historial de conversión tiene antecedentes consolidados.
\citeA{gonzalezflores2025relevance} documentan un caso en el que un esquema tradicional de lead scoring, basado en asignación manual de valores por parte de marketing y ventas, fue reemplazado por un modelo predictivo seleccionado tras comparar quince algoritmos de clasificación, logrando mejoras sustanciales en la capacidad de distinguir leads de alta y baja calidad. Su diagnóstico de que los métodos manuales tienden a ser arbitrarios y menos precisos coincide con el del proceso comercial analizado en este trabajo.

En la literatura de \emph{lead scoring} se han comparado algoritmos como la
regresión logística y los métodos de ensamble \citeA{nygard2020automating}, con
resultados que dependen de las características del conjunto de datos, la
definición del objetivo y las métricas de evaluación empleadas.
\citeA{tan2024leveraging} obtienen resultados equivalentes en el sector educativo,
apoyándose en variables de comportamiento digital como la duración de visita a las
páginas.

La revisión de \citeA{wu2023state} sobre el estado de los modelos de calificación y
su impacto en el desempeño comercial matiza estos resultados: la ganancia
documentada depende menos del algoritmo seleccionado que de la calidad de las
variables disponibles y de la integración efectiva del modelo en el flujo de
trabajo del equipo de ventas. Esta observación es consistente con
\citeA{ahearne2007sales}, quienes muestran que el impacto de los sistemas CRM sobre
la efectividad comercial se materializa únicamente cuando los representantes
incorporan la herramienta en su rutina.

\section{Capacidad comercial y priorización bajo restricción}

\citeA{sabnis2013blackhole} documentan el fenómeno del \emph{agujero negro de los
leads}: la fracción de oportunidades generadas por mercadeo que no son perseguidos debido a la capacidad limitada del equipo comercial. Su aporte al
presente trabajo es conceptual: establece que el problema no es de generación sino
de asignación, lo que justifica formular la evaluación en términos de conversiones
capturadas bajo una restricción de capacidad, y no en términos de exactitud
agregada.

Los estudios sobre el tiempo de respuesta a leads web muestran que las
probabilidades de contacto y calificación disminuyen a medida que aumenta el tiempo
transcurrido desde la solicitud. \citeA{oldroyd2011short} encuentran que intentar
contactar a un prospecto dentro de la primera hora se asocia con una mayor
probabilidad de calificación que hacerlo después de ese lapso. En consecuencia, bajo
una capacidad limitada de atención comercial, la oportunidad de respuesta debe
considerarse junto con la probabilidad de conversión al definir el orden de
priorización.

\citeA{ohiomah2019lead} aportan evidencia empírica de que el uso de sistemas de
gestión de leads se asocia con un mejor desempeño en ventas internas, mediante una
mayor intensidad de seguimiento de leads, una venta más adaptativa y el
fortalecimiento de competencias técnicas y de venta. Este hallazgo sustenta la
decisión de entregar la priorización como un artefacto operativo integrado al
proceso de gestión de leads, y no únicamente como un reporte analítico
independiente.

\section{Desbalance de clases y evaluación orientada al beneficio}

En problemas de clasificación binaria con clase minoritaria muy escasa,
\citeA{verbeke2012new} desarrollan un enfoque explícitamente orientado al
beneficio, basado en el beneficio máximo potencial de intervenir sobre una fracción
óptima de clientes, y muestran que seleccionar modelos mediante métricas
exclusivamente estadísticas puede conducir a decisiones subóptimas.
\citeA{burez2009handling} examinan el tratamiento del desbalance de clases
empleando AUC y \emph{lift}, y encuentran que el submuestreo avanzado CUBE no
mejora el desempeño predictivo, lo que indica que las estrategias de remuestreo no
garantizan mejoras en métricas de discriminación u ordenamiento. De forma
consistente, en el presente trabajo SMOTE no mejoró las métricas de ordenamiento y
deterioró la calibración; este resultado debe interpretarse como un hallazgo
específico de los datos y modelos analizados, no como una consecuencia general del
remuestreo.

En una línea complementaria, \citeA{devriendt2018literature} revisan el modelado de
elevación (\emph{uplift modeling}), cuyo objetivo es estimar el efecto incremental
de una intervención sobre un resultado, en lugar de predecir únicamente la
probabilidad de dicho resultado. Este enfoque quedó fuera del alcance del presente
trabajo, pues requiere datos que permitan identificar el efecto de una acción
comercial frente a una condición de comparación —idealmente mediante un diseño
experimental aleatorizado—. No obstante, constituye una extensión pertinente si se
dispone de intervenciones comerciales observables y datos adecuados para estimar sus
efectos causales.

\section{Segmentación de clientes}

La segmentación de clientes mediante técnicas de agrupamiento constituye una práctica consolidada en investigación de mercados y analítica comercial. \cite{kaufman1990finding} presentan fundamentos y métodos clásicos de análisis de conglomerados. Estas técnicas pueden emplearse, entre otros fines, para apoyar la definición de territorios comerciales, la asignación de carteras y la identificación de perfiles de clientes, siempre que se valide empíricamente la utilidad operativa de los segmentos obtenidos.

En la base de datos utilizada en este trabajo, varios atributos relevantes, por ejemplo, motivación de compra, estado civil y tipología de unidad de interés, son categóricos. Esta característica exige precaución al utilizar K-Means estándar, ya que este algoritmo se basa en centroides y distancias euclídeas, y por tanto está definido de forma natural para variables numéricas. Para datos categóricos o mixtos, pueden emplearse codificaciones adecuadas junto con criterios de escalamiento y ponderación, o bien métodos específicamente diseñados para este tipo de información, como K-prototypes o KAMILA. La estrategia adoptada en este trabajo debe especificar la codificación aplicada, la métrica de distancia y el criterio usado para fijar los pesos, a fin de evitar que el número de categorías o la escala de un bloque de variables domine la segmentación.

\section{Modelado predictivo en el sector inmobiliario}

La investigación aplicada al sector inmobiliario ha documentado que la combinación
de juicio experto con modelos analíticos habilita la identificación de
oportunidades no evidentes \cite{palmer2025bridging}. La complejidad
particular del sector, tales como ciclos de decisión largos, desembolso elevado y participación
de varias personas en la decisión, reduce la eficacia de los enfoques basados
exclusivamente en reglas estáticas.


\section{Aporte diferencial del presente trabajo}

Respecto a los antecedentes revisados, este trabajo aporta cuatro elementos:

\begin{enumerate}
    \item \textbf{Tratamiento de la ausencia de datos como señal.} La literatura
    consultada asume habitualmente conjuntos de datos con completitud razonable. El
    contexto de este trabajo: un CRM donde el perfil se completa progresivamente
    durante la calificación, obliga a un tratamiento distinto, en el que la propia
    estructura de ausencia se incorpora como información predictiva en lugar de
    eliminarse mediante imputación.
    \item \textbf{Evaluación contra una distribución de referencia explícita.} En
    lugar de reportar únicamente métricas de discriminación, el desempeño se
    contrasta contra la distribución completa del resultado que produciría la
    ausencia de priorización, obtenida por simulación y verificada analíticamente
    contra la distribución hipergeométrica exacta.
    \item \textbf{Restricción deliberada de generalización.} La exclusión de la
    identidad del proyecto del conjunto de predictores, con cuantificación explícita
    del costo estadístico que implica, responde al requisito de aplicabilidad a
    proyectos residenciales futuros.
    \item \textbf{Integración del análisis de procesos con el modelo predictivo.} El
    modelado BPMN del proceso comercial no se limita a documentar el estado actual:
    determina el punto de inserción de la intervención analítica y define el momento
    en que la predicción debe emitirse, lo que a su vez fija qué variables pueden
    considerarse disponibles y cuáles constituyen fuga de información.
    \item \textbf{Contextualización en el mercado inmobiliario guatemalteco.}
    El trabajo aporta evidencia empírica obtenida en el contexto de empresas
    desarrolladoras de vivienda en la Ciudad de Guatemala, a partir de datos de
    proyectos residenciales locales. Esta contextualización permite evaluar el
    comportamiento del enfoque de priorización bajo las características propias
    del entorno comercial analizado, en lugar de asumir directamente que los
    resultados obtenidos en otros contextos son transferibles al mercado
    guatemalteco.
\end{enumerate}
